\scanpage{62}

% manuscript line(s) 1
Je vois, du côté $E$, la condition que
% manuscript line(s) 2
\[
Q_{\deleted{\unsure{x}}T,N}(x)\overset{\mathrm{def}}=x^2\deleted{\unsure{+}}\unsure{-|+}xT+N\qquad(x\in\Gamma P_1)
\]
% manuscript line(s) 3
à valeurs dans $L^{\otimes2}$, ou \unsure{mieux|encore}
% manuscript line(s) 4
\[
Q_{T,N}(x)=x^2-\psi(x)xT+\psi(x)^2N\qquad(x\in\Gamma\underline E)
\]
% manuscript line(s) 5–6
à valeurs dans $L^{\otimes2}$. Faisant $x=T$, on trouve la condition familière (compte tenu de $\psi(T)=2$)
% manuscript line(s) 7
\[
4N-T^2\in\Gamma L^{\otimes2}\quad\subset\Gamma\mathrm{Sym}^2 E
\]
% manuscript line(s) 8–10
mais elle ne suffit pas. En termes d'un scindage $E\simeq L\oplus\underline O$ par choix d'un $x_0\in\Gamma P_1$, on aura
% manuscript line(s) 10
\[
\mathrm{Sym}^2E\simeq\underline L^2\oplus\underset{\substack{\shortparallel\\x_0\cdot L}}{\underline L}\oplus\underset{\substack{\phantom{\shortparallel}\\=x_0^2\underline O_X}}{\underline O}
\]
% manuscript line(s) 11–13
\[
\begin{aligned}&\deleted{\unsure{y}}\quad x_0=(0,1)\deleted{\unsure{x}}\\ &T=(t,2)=2x_0+\underset{(t,0)}t,\quad t\in\Gamma L,\\ &N=\underset{\substack{\in\\\Gamma L^{\otimes2}}}c-x_0\underset{\substack{\in\\\Gamma L}}b+x_0^2\quad(\text{car }\psi_2(N)=1)\end{aligned}\qquad\begin{cases}t\in\Gamma L\\ b\in\Gamma L\\c\in\Gamma L^{\otimes2}\end{cases}
\]
% manuscript line(s) 14
Soit alors $x=(u,1)\inserted{=u+x_0}$, $u\in L$, on aura
% manuscript line(s) 15–17
\[
\begin{aligned}Q_{T,N}(x_0+u)=Q_{T,N}(x)&=(x_0+u)^2-(x_0+\underset{\substack{\in\\\Gamma L}}u)(2x_0+\underset{\substack{\in\\\Gamma L}}t)+\underset{\substack{\in\\\Gamma L^{\unsure{2}}}}c-\underset{\substack{\in\\x_0\otimes L}}{bx_0}+x_0^2\\ &\deleted{\unsure{=2x_0u+x_0^2+u^2-(x_0+u)t}}\\ &=\underset{\substack{\in\\L^{\otimes2}}}{u^2}-\underset{\substack{\in\\L^{\otimes2}}}{ut}+\deleted{\unsure{x_0^2+}}(-bx_0+c-x_0t)\end{aligned}
\]
% manuscript line(s) 18
et la condition est que
% manuscript line(s) 19–20
\[
\begin{aligned}&\deleted{x_0^2-x_0(x_0}\\&\underbrace{x_0(\deleted{\unsure{x_0+}}-b-t)}_{\in\Gamma(E\otimes L)}+\underset{\substack{\in\\\Gamma(L^{\otimes2})\subset\Gamma(E\otimes L)}}c\in\Gamma L^{\otimes2}\end{aligned}
\]
% manuscript line(s) 21
Cela signifie $\boxed{t=-b}$ i.e. la condition est \deleted{que}
% manuscript line(s) 21–22
% slanted
% manuscript line(s) 21–22
\margin{Ceci \deleted{signifie} \inserted{implique} que sur le \unsure{composant} du scindage de l'ext. $\mathcal A$ de $\underline O$ par $L$, $b$ est le \unsure{même} (le \unsure{terme en} \unsure{bon} au signe}
% manuscript line(s) 22
\[
\boxed{N=x_0^{\unsure{2}}-\underset{\substack{\in\\\Gamma L}}b x_0+\underset{\substack{\in\\\Gamma L^{\otimes2}}}c\quad\Longrightarrow\quad T=2x_0\deleted{\unsure{+}}-b}
\]
% manuscript line(s) 22–23
% slanted
% manuscript line(s) 22–23
\margin{d'où $\boxed{\unsure{\Theta|D|\delta}=T^2-4N=b^2-4c}$}
% manuscript line(s) 23
d'où $\boxed{Q(x_0+u)=u^2+\deleted{\unsure{t}}bu+c}$
% archive sheet number: 61
